\section{Metodologia}
\label{sec:background}

A metodologia adotada neste trabalho foi estruturada de forma a avaliar, de maneira sistemática e controlada, a eficácia da arquitetura proposta para segurança em Modelos de Linguagem de Grande Escala (LLMs). O processo metodológico integra a geração de cenários de ataque, a classificação de riscos, a execução de testes em múltiplos idiomas e a análise quantitativa dos resultados, permitindo uma investigação abrangente das capacidades e limitações do sistema.

\subsection{Estratégia Experimental}

O experimento foi concebido como um \textit{ataque exploratório contínuo}, cujo objetivo é simular o comportamento de um agente malicioso persistente. Essa abordagem consiste na geração automática de entradas projetadas para induzir o LLM a violar políticas de segurança, produzir conteúdo nocivo, revelar informações sensíveis ou demonstrar vieses linguísticos. A estratégia busca imitar situações realistas, considerando que agentes adversários frequentemente exploram múltiplas técnicas, idiomas e estruturas linguísticas para burlar mecanismos de defesa.

A escolha por uma metodologia exploratória contínua permitiu avaliar não apenas casos isolados, mas também a capacidade do sistema de manter níveis consistentes de segurança diante de ataques sucessivos.

\subsection{Geração Multilíngue de Cenários}

Um dos diferenciais desta metodologia é o uso de cenários de ataque em múltiplos idiomas, incluindo português, inglês e espanhol. A literatura demonstra que diversos ataques de \textit{prompt injection} e \textit{jailbreak} apresentam variações linguísticas significativas, sendo possível contornar filtros monolíngues por meio de:

\begin{itemize}
    \item mudanças de idioma;
    \item alternância entre códigos linguísticos (\textit{code-switching});
    \item tradução contextual com semântica preservada;
    \item inserção de ambiguidades culturais específicas.
\end{itemize}

Ao incluir múltiplos idiomas, o experimento adotou uma perspectiva mais realista e alinhada ao contexto global de uso dos LLMs, ampliando a cobertura metodológica e expondo limitações que dificilmente seriam observadas em cenários monolíngues. A literatura existente ainda não explora de forma ampla essa dimensão, reforçando a originalidade da abordagem.

\subsection{Construção dos Casos de Teste}

Foram desenvolvidos 30 casos de teste distribuídos entre quatro grupos principais:

\begin{itemize}
    \item \textbf{Ataques explícitos}: tentativas de \textit{prompt injection}, indução maliciosa e manipulação direta do comportamento do modelo;
    \item \textbf{Detecção de PII}: solicitações envolvendo CPF, e-mail, telefone, cartão de crédito e número de seguridade social;
    \item \textbf{Vieses linguísticos}: categorias relacionadas a gênero, raça, idade, corpo, orientação, localização geográfica e educação;
    \item \textbf{Alucinações}: prompts projetados para induzir respostas factualmente incorretas ou inconsistentes.
\end{itemize}

Cada caso de teste foi submetido à arquitetura completa, contemplando sanitização, verificação de entrada, processamento pelo LLM e auditoria na saída.

\subsection{Métricas de Avaliação}

As seguintes métricas foram utilizadas para análise quantitativa dos resultados:

\begin{itemize}
    \item \textbf{Precisão} (\textit{Precision}) — mede a proporção de classificações corretas entre as respostas identificadas como risco;
    \item \textbf{Revocação} (\textit{Recall}) — indica a capacidade do sistema de detectar todos os riscos reais;
    \item \textbf{F1-score} — combinação harmônica entre precisão e recall, útil em cenários com desequilíbrio entre classes;
    \item \textbf{Acurácia} — porcentagem geral de classificações corretas;
    \item \textbf{Taxa de sucesso} — percentual de casos em que a ameaça foi devidamente identificada e bloqueada;
    \item \textbf{Tempo de resposta} — medição do custo computacional associado a cada módulo.
\end{itemize}

Essas métricas permitem identificar tanto o desempenho geral quanto a sensibilidade do sistema a categorias específicas de ameaças.

\subsection{Pipeline de Execução}

O pipeline experimental segue cinco etapas principais:

\begin{enumerate}
    \item Geração automática ou manual dos prompts adversariais.
    \item Execução do fluxo completo da arquitetura (sanitização, verificação, LLM e auditoria de saída).
    \item Registro estruturado dos resultados, incluindo classificação, métricas e tempos de resposta.
    \item Cálculo das métricas por categoria e comparação entre classes.
    \item Análise dos dados, interpretação dos padrões e identificação de limitações.
\end{enumerate}

O uso de um pipeline padronizado permite a reprodutibilidade dos experimentos e facilita auditorias futuras.

\subsection{Resumo}

De forma geral, a metodologia adotada incorpora elementos exploratórios, linguísticos, quantitativos e comportamentais, proporcionando uma avaliação abrangente da robustez da arquitetura. A abordagem multilíngue, a simulação contínua de ataques e o uso de métricas rigorosas ampliam significativamente a validade dos resultados apresentados na Seção~\ref{sec:resultados_finais}.
